\chapter{Superficies en el espacio}

\begin{definicion}
    Una \textbf{superficie} en $\bb{R}^3$ es un subconjunto $S\subset \bb{R}^3$ no vacío tal que $\forall p \in S$ existe\footnote{cada uno de estos elementos deberían tener subíndice $p$ porque dependen del punto pero por simplificar la notación no lo pondremos}
    \begin{enumerate}
        \item $U\subset \bb{R}^2$ abierto.
        \item un entorno abierto $V$ de $p$ en $S$.
        \item una aplicación $X:U \to \bb{R}^3$ diferenciable.
    \end{enumerate}
    que cumplen:
    \begin{enumerate}
        \item $X(U)=V$
        \item $X:U\to V$ es un homeomorfismo.
        \item $(dX)_q : \bb{R}^2 \to \bb{R}^3$ inyectiva $\forall q=(u,v) \in U$.
    \end{enumerate}
    Normalmente notaremos
    \Func{X}{U}{V}{(u,v)}{(x(u,v), y(u,v), z(u,v))}
    donde $x,y,z:U\to \bb{R}$ son diferenciables. Además:
    \begin{itemize}  
            \item A $(u,v)$ lo llamaremos \textbf{coordenadas} de $p$.
            \item A $X$ lo llamaremos \textbf{parametrización} (local).
            \item A $V$ lo llamaremos \textbf{entorno parametrizado}.
            \item A $U$ lo llamaremos \textbf{dominio de la parametrización}.
        \end{itemize}
\end{definicion}

\begin{observacion}
    Con esta definición no puede haber superficies con autointersecciones (ya que en ese caso $X$ no sería biyectivo)\footnote{Realmente sí existen pero no las vamos a estudiar este año.}.
\end{observacion}

\begin{notacion}
    Hemos escrito $(dX)_q$, lo cual hasta ahora habremos notado como $(DX)(q,v)$ que es la derivada direccional en la dirección del vector $v$ del punto $q$. En este caso no especificamos la dirección $v$ porque se considera en todas las direcciones (sticker de shrek).
\end{notacion}

\begin{observacion}
    Otra condición equivalente a que $(dX)_q$ sea inyectiva para todo $q\in U$ sería que para todo $q\in U$ se tenga que
    \begin{gather*}
        \frac{\partial X}{\partial u}(q),\ \ \ \frac{\partial X}{\partial v}(q)\ \ \ \text{ son linealmente independientes}.
    \end{gather*}
    Lo cual además sería equivalente a decir que 
    \begin{gather*}
        \left(\frac{\partial X}{\partial u}(q) \wedge \frac{\partial X}{\partial v}(q)\right) \neq 0\ \ \ \ \forall q\in U
    \end{gather*}
\end{observacion}

\begin{ejemplo}\
    \begin{enumerate}
        \item Cualquier plano $P\subset \bb{R}^3$ se puede escribir como
        \begin{gather*}
            P = \{(x,y,z)\in \bb{R}^3 : ax+by+cz+d=0\}\ \ \ (a,b,c)\neq (0,0,0),\ \ d\in \bb{R}
        \end{gather*}
        Supongamos que $c\neq 0$, entonces 
        \begin{gather*}
            P=\{(x,y,z)\in \bb{R}^3 : z=Ax+By+C\}\ \ \ A,B,C \in \bb{R}
        \end{gather*}
        Veamos que $P$ es una superficie. Para ello, para cada $p\in P$ tenemos que encontrar $U\subset \bb{R}^2$, $V\subset \bb{R}^3$ y $X:U\to V$ que cumplan las condiciones de la definición. Sea $p\in P$ elegimos
        \begin{gather*}
            U=\bb{R}^2, \ \ \ \ V=P
        \end{gather*}
        \Func{X}{U}{\bb{R}^3}{(u,v)}{(u,v,Au+Bv+C)}
        que vemos que no dependen del punto $p$ (ya que estamos en un caso muy sencillo). Además, claramente $X$ es diferenciable. Veamos que se verifican las condiciones para ser una superficie con estos elementos seleccionados:
        \begin{enumerate}
            \item $X(U)=X(\bb{R}^2) = \{(x,y,z)\in \bb{R}^3 : (x,y)\in U, z=Ax+By+C\} = P = V$.
            
            \item Veamos que $X: U=\bb{R}^2 \to V=P$ es un homeomorfismo. Es claramente continua y sobreyectiva, nos falta ver que es inyectiva, lo cual se tiene ya mirando las dos primeras componentes. Podemos considerar su inversa $\pi: P \to \bb{R}^2$ dada por 
            \begin{gather*}
                \pi(x,y,z) = (x,y) \ \ \ \ \forall (x,y,z)\in P
            \end{gather*}
            que es claramente continua. Confirmamos que $X$ es un homeomorfismo.

            \item Se tiene para todo $(u,v)\in U$:
            \begin{gather*}
                X_u(u,v) = (1,0,A)\\
                X_v(u,v) = (0,1,B)
            \end{gather*}
            y es fácil ver que son linealmente independientes (mirando las 2 primeras componentes de cada derivada parcial).
        \end{enumerate}
        \item Veamos que los grafos de funciones son superficies. Para ello consideramos $O\subset \bb{R}^2$, $f:O\to \bb{R}$ diferenciable y tendremos
        \begin{gather*}
            Grafo(f) = \{(x,y,z)\in \bb{R}^3:(x,y)\in O,\ \ z=f(x,y)\}
        \end{gather*}
        Tomamos $p\in Grafo(f)$ y elegimos
        \begin{gather*}
            U=O,\ \ \ \ V=Grafo(f)
        \end{gather*}
        \Func{X}{U}{\bb{R}^3}{(u,v)}{(u,v,f(u,v))}
        donde $X$ es diferenciable por serlo la identidad y $f$. Veamos que cumple las condiciones para ser una superficie:
        \begin{enumerate}
            \item $X(O) = \{(u,v,z)\in \bb{R}^3 : (u,v)\in O,\ \ z=f(u,v)\} = Grafo(f) = V$.
            \item En topología vimos que todo grafo es homeomorfo a su dominio (con la proyección como inversa) por lo que $X$ es un homeomorfismo.
            \item $X_u(u,v) = (1,0,f_u)(u,v)$ y $X_v(u,v)=(0,1,f_v)(u,v)$ que con las dos primeras componentes ya vemos que son linealmente independientes.
        \end{enumerate}
        En este caso vemos de nuevo que estos elementos no dependen del punto $p$ escogido.

        % \item Consideramos el siguiente problema:
        % \begin{gather*}
        %     \begin{array}{l}
        %         X(u,v) = (x,y,z)\\
        %         X(q) = p
        %     \end{array}
        % \end{gather*}
        % \begin{itemize}  
        %     \item A $(u,v)$ lo llamaremos \textbf{coordenadas} de $p$.
        %     \item A $X$ lo llamaremos \textbf{parametrización} (local).
        %     \item A $V$ lo llamaremos \textbf{entorno parametrizado}.
        %     \item A $U$ lo llamaremos \textbf{dominio de la parametrización}.
        % \end{itemize}

        \item Los abiertos de una superficie son una superficie
        
        \item Sea $S\subset \bb{R}^3$ no vacío y supongamos que $S_i\subset S$ es una superficie abierta para cada $i\in I$ tal que
        \begin{gather*}
            S=\bigcup\limits_{i\in I}S_i
        \end{gather*} 
        Entonces $S$ es una superficie.\\

        Para verlo tomamos $p\in S$ por lo que existirá un $i_0\in I$ tal que $p\in S_{i_0}$ superficie. Por tanto tenemos que 
        \begin{gather*}
            \exists U_0\subset \bb{R}^2 \text{ abierto },\ \ \exists V_0 \text{ entorno abierto de $p$ en $S_{i_0}$},\ \ \exists X:U_0\to \bb{R}^3\text{ diferenciable}
        \end{gather*}
        cumpliendo las 3 propiedades de una superficie. Elegimos $V=V_0\subset S$ que es un entorno abierto de $p$ en $S_{i_0}\subset$, que es abierto en $S$. Tomamos $U=U_0$ que es abierto en $\bb{R}$ y tomamos $X_0=X:U=U_0 \to \bb{R}^3$ y se deja como ejercicio ver que se verifican las propiedades de una superficie.

        \item Sea $S\subset \bb{R}^3$ una superficie, $O,O'\subset \bb{R}^3$ abiertos con $S\subset O$ y consideramos $\phi:O\to O'$ un difeomorfismo. Entonces $S'=\phi(S)$ es otra superficie en $\bb{R}^3$.\\
        
        Tomamos $p'\in S'$ y tenemos que $\exists!p\in S$ tal que $\phi(p)=p'$. Como $S$ es superficie, asociados a $p$ existirán $U,V$ y $X$ verificando las propiedades de una superficie. Tomaremos $U'=U$, $V'=\phi(V)$ y $X'=\phi \circ X$. Veamos que cumplen lo que buscamos:
        \begin{enumerate}
            \item $X'(U)=(\phi \circ X)(U)=\phi(X(U)) = \phi(V)=V'$.
            \item Veamos que $X':U'\to V'$. Para ello estamos en la siguiente situación
            \begin{gather*}
                \xymatrix{
                    \pi \circ X: U' \ar[d]  \ar[r] & V'\\
                    V \ar[ur]
                }
            \end{gather*}

            \item Sea $q'\in U'=U$, entonces tenemos que 
            \begin{gather*}
                (dX')_q = d(\phi \circ X)_{q'}
            \end{gather*}
        \end{enumerate}

        \item Consideramos las esferas $S^2(c, R)$ con $c\in \bb{R}^3$, $R>0$. Supongamos $c=0$ y $R=1$ (sin pérdida de generalidad). Entonces tenemos que $S^2(c, R) = S^2$. Sea $N=(0,0,1)\in S^2$ y elijo 
        \begin{gather*}
            U=\{(u,v)\in \bb{R}^2 : u^2 + v^2 <1\}
        \end{gather*}
        que es abierto en $\bb{R}^2$. Elegimos $V$ como el hemisferio abierto de centro $N$ en $S^2$, es decir,
        \begin{gather*}
            V = H_N^+ = \{(x,y,z)\in S^2 : z>0\}
        \end{gather*}
        y tenemos que $N\in H_N^+$, por lo que $V$ es entorno abierto de $N$. Elegimos además la aplicación dada por 
        \Func{X}{U}{\bb{R}^3}{(u,v)}{(u,v,\sqrt{1-u^2 - v^2})}
        que es diferenciable. Tenemos entonces
        \begin{enumerate}
            \item \begin{align*}
                X(U)&=\{(x,y,z)\in \bb{R}^3 : (x,y)\in U, z=\sqrt{1-u^2-v^2}\} =\\
                &= \{(x,y,z)\in S^2 : z>0\} = H_N^+ = V
            \end{align*}

            \item Vemos que $X:U\to V$ y además 
            \Func{X^{-1}}{V}{U}{(x,y,z)}{(x,y)}
            \item Tenemos 
            \begin{gather*}
                X_u(u,v) = (1,0,\ast)(u,v)\\
                X_v(u,v) = (0,1,\ast)(u,v)
            \end{gather*}
            Por lo que $X_u$ y $X_v$ son linealmente independientes en cada $(u,v)\in U$.
        \end{enumerate}
        Con esto habríamos probado que se cumplen las condiciones para $N$ pero tenemos que ver que se cumplen para todo $p\in S^2$. Para ello consultamos ejemplos anteriores y vemos que dado $p\in S^2$ existe un giro $\phi_p : \bb{R}^3 \to \bb{R}^3$ difeomorfismo tal que $\phi_p(N)=p$. Entonces $\phi_p(H_N^+) = H_p^+$ por lo que $H_p^+$ es una superficie.\\

        Sea $S=S^2(c, R)$, entonces existen una traslación $T$ y una homotecia $H$ de $\bb{R}^3$ tales que 
        \begin{gather*}
            (H\circ T)(S^2) = S^2(c,R)
        \end{gather*}
        por lo que $S^2(c, R)$ es una superficie. Con este razonamiento llegamos a que las elipsoides también son superficies.

        \item Superficies en forma implícita. Sea $O\subset \bb{R}^3$ abierto $f:O\to \bb{R}$ diferenciable. Tenemos que $a\in \bb{R}$ se llama \textbf{valor regular} de $f$ cuando
        \begin{gather*}
            (df)_p = (\nabla f)_p \neq 0 \ \ \ \ \forall p\in f^{-1}(\{a\})
        \end{gather*}
        Consideramos $S=f^{-1}(\{a\})\subset O \subset \bb{R}^3$ con $a$ un valor regular y queremos ver si $S$ es o no una superficie. Recordamos que el Teorema de la Función implícita dice que para cada punto $p_0 = (x_0, y_0, z_0)\in S$ si se tiene que $(\nabla f)_{p_0}\neq 0$, entonces suponemos sin pérdida de generalidad que $f_x(p_0)\neq 0$. Tenemos que existe un entorno abierto $U$ de $(x_0,y_0)$ en $\bb{R}^2$, un entorno abierto $V$ de $z_0$ en $\bb{R}$ y una aplicación $g:U\to V\subset \bb{R}$ diferenciable tales que 
        \begin{gather*}
            S \cap (U\times V) = Grafo(g)
        \end{gather*}
        Por tanto $S$ es unión de grafos abiertos en $S$.

        \item Sea $A\subset M_4(\bb{R})$ tal que $A=A^t$. Sea
        \begin{gather*}
            S = \{p\in \bb{R}^3 : (1,p^t)A \begin{pmatrix} 1\\p\end{pmatrix} = 0\} 
        \end{gather*}
        que será la cuádrica en $\bb{R}^3$ determinada por la matriz\footnote{es la matriz ampliada, pero aquí no escribimos eso} $A$. Queremos ver ahora entonces si las cuádricas de $\bb{R}^3$ son superficies.\\

        Si definimos $f:O=\bb{R}^3 \to \bb{R}$ dada por 
        \begin{gather*}
            f(p) =  (1,p^t)A \begin{pmatrix} 1\\p\end{pmatrix}
        \end{gather*}
        tendremos que es diferenciable y además con esta definición se tiene que $S=f^{-1}(\{0\})$. Nos preguntamos ahora si $0$ es un valor regular de $f$ (ya que en dicho caso podremos aplicar el apartado anterior). Aplicando la definición de valor regular tomamos $p\in f^{-1}(\{0\})=S$ y consideramos $(df)_p:\bb{R}^3 \to \bb{R}$. Tomamos $v\in \bb{R}^3$ y se tiene 
        \begin{gather*}
            (df)_p(v) = (0,v^t)A\begin{pmatrix} 1\\p\end{pmatrix} + (1,p^t)A\begin{pmatrix} 0\\v\end{pmatrix} = 2(1,p^t)A\begin{pmatrix} 0\\v\end{pmatrix}
        \end{gather*}
        Supongamos que $0$ no es valor regular. En ese caso existirá un $p\in f^{-1}(\{0\})$ tal que $(df)_p = 0$ lo que nos diría que 
        \begin{gather*}
            (df)_p(v) = 0 \ \ \ \ \forall v \in \bb{R}^3
        \end{gather*}
        o equivalentemente
        \begin{gather*}
            0 = (df)_p(v) = 2(1,p^t)A \begin{pmatrix} 0\\v\end{pmatrix} \ \ \ \ \forall v \in \bb{R}^3 \Rightarrow (1,p^t)A = (\lambda, 0)
        \end{gather*}
        Si lo unimos con lo anterior tenemos que 
        \begin{gather*}
            0 = f(p) = (1,p^t)A \begin{pmatrix} 0\\v\end{pmatrix} = (\lambda, 0)\begin{pmatrix} 1\\p\end{pmatrix} = \lambda
        \end{gather*}
        Si $0$ no es valor regular, entonces existe un $p\in S$ tal que $(1,p^t)A = 0$. El contrarrecíproco de esta afirmación sería que si para todo $p\in S$ se tiene que $(1,p^t)A\neq 0$, entonces se tiene que $0$ es regular para $f$. En este caso tendríamos que o bien $S$ es vacío o bien $S$ es una superficie.\\
        
        Si la cuádrica $S$ no tiene punto singulares (los tiene en el infinito) y no es vacía, entonces es una superficie de $\bb{R}^3$. Esto nos dice que los elipsoides, hiperboloides, paraboloides, cilindros y planos son superficies. Nos quedarían fuera los conos y los planos secantes, que en principio no podremos concluir nada con el ejemplo anterior.

        \item Consideramos $S^1((0,a,0), r)$ con $0<r<a$. Definimos el siguiente conjunto:
        \begin{gather*}
            S = \{(x,y,z)\in \bb{R}^3 : (x,y,z) \text{ se obtiene de algún punto de }S^1((0,a,0), r)\\
            \text{ rotando alrededor del eje }z\}
        \end{gather*}
        que llamaremos \textbf{toro de revolución}\footnote{revolución no porque sea muy rebelde o porque quiera derrocar al gobierno.}. Tomamos $(x,y,z)\in S$ por lo que se habrá obtenido de un punto de $S^1((0,a,0), r)\owns (0,a+r\cos(u), r\sen(u))$. Las rotaciones alrededor del eje $z$ serán producto de multiplicar por la siguiente matriz:
        \begin{gather*}
            \begin{pmatrix}
                \cos(v) & \sen(v) & 0\\
                -\sen(v) & \cos(v) & 0\\
                0 & 0 & 1
            \end{pmatrix}
        \end{gather*}
        Por tanto, al aplicar el giro sobre el eje $z$ tendremos:
        \begin{align*}
            (x,y,z) &= (0,a+r\cos(u), r\sen(u)) \begin{pmatrix}
                \cos(v) & \sen(v) & 0\\
                -\sen(v) & \cos(v) & 0\\
                0 & 0 & 1
            \end{pmatrix}
            =\\ &=((a+r\cos(u))\sen(v), (a+r\cos(u))\cos(v), r\sen(u)) \ \ \ \ \forall u,v\in \bb{R}
        \end{align*}
        Esto sería una construcción discreta de la parametrización $X(u,v)$. Si estudiamos la condición que nos ha salido antes nos quedará
        \begin{gather*}
            x^2 + y^2 = (a+r\cos(u))^2 \Rightarrow \sqrt{x^2 + y^2} = a+r\cos(u) \Rightarrow \sqrt{x^2 + y^2} -a = r\cos(u)
        \end{gather*}
        Sabemos que $z=r\sen(u)$ por lo que juntando estas ecuaciones nos queda
        \begin{gather*}
            (\sqrt{x^2 + y^2}-a)^2 + z^2 = r^2
        \end{gather*}
        Entonces $S\subset \{(x,y,z)\in \bb{R}^3 : (\sqrt{x^2 + y^2}-a)^2 + z^2 = r^2\}$ y queremos buscar la igualdad. Para ello tomamos $(x,y,z)$ de este conjunto, y tendremos
        \begin{gather*}
            \frac{(\sqrt{x^2 + y^2}-a)^2 + z^2}{r^2} = 1
        \end{gather*}
        por lo que para cierto $u\in \bb{R}$ se tendrá que
        \begin{gather*}
            \left\{
            \begin{array}{l}
                \sqrt{x^2 + y^2} -a = r\cos(u) \Rightarrow \sqrt{x^2 + y^2} = a + r\cos(u) \Rightarrow x^2 + y^2 = (a+r\cos(u))^2\\
                z = r\sen(u) 
            \end{array}
            \right.
        \end{gather*}
        de la primera ecuación llegamos a que 
        \begin{gather*}
            \left\{
            \begin{array}{l}
                x = (a+r\cos(u))\cos(v)\\
                y = (a+r\cos(u))\sen(v)
            \end{array}
            \right. \ \ \ \ v\in \bb{R}
        \end{gather*}
        Definimos $f(x,y,z) = (\sqrt{x^2 + y^2} -a)^2 + z^2$ que será diferenciable en $O=\bb{R}^3 \setminus\{\text{eje }z\}$. Nos preguntamos ahora si $r^2$ es un valor regular de $f$. Para ello tomamos $(x,y,z)\in f^{-1}(\{r^2\})=S$ y estudiamos su gradiente:
        \begin{gather*}
            (\nabla f)_{(x,y,z)} = \left(\frac{2(\sqrt{x^2 + y^2} -a)x}{\sqrt{x^2 + y^2}}, \frac{2(\sqrt{x^2 + y^2} -a)y}{\sqrt{x^2 + y^2}}, 2z \right) 
        \end{gather*}
        Si $r^2$ no es regular para $f$, entonces existe un $(x,y,z)\in S$ para el cual $(\nabla f)_{(x,y,z)} = (0,0,0)$. Si escribimos esto tendremos
        \begin{gather*}
            \left\{
            \begin{array}{l}
                2(\sqrt{x^2 + y^2} -a)x = 0\\
                2(\sqrt{x^2 + y^2} -a)y = 0\\
                2z = 0 \Rightarrow z=0
            \end{array}
            \right.
        \end{gather*} 
        Si distinguimos casos tenemos 
        \begin{itemize}
            \item Si $\sqrt{x^2 + y^2}-a \neq 0$, entonces $x=y=z=0$ pero entonces $(x,y,z)\notin O$.
            \item Si $\sqrt{x^2 + y^2}-a = 0$, entonces $f(x,y,z) = 0 \neq r^2$
        \end{itemize}
        En cualquier caso llegamos a una clara contradicción por lo que $r^2$ es un valor regular. Por el ejemplo 7 sabemos que $S$ o bien es vacío o bien es una superficie. Como la circunferencia inicial estaba en $S$ sabemos que no es vacío por lo que concluimos finalmente que el toro de revolución es una superficie.
    \end{enumerate}
\end{ejemplo}

\begin{ejercicio}
    El cono $\cc{C}$ dado por 
    \begin{gather*}
        \cc{C} = \{(x,y,z)\in \bb{R}^3 : x^2 + y^2 = z^2\}
    \end{gather*}
    no es una superficie.\\

    Si $\cc{C}$ fuera una superficie entonces existirían $V$ entorno del $(0,0,0)$ y un abierto $U\subset \bb{R}^2$ conexo tales que existe un homeomorfismo $X$ entre ellos. Sin embargo cualquier entorno abierto $V$ del $(0,0,0)$ verificará que $V\setminus\{(0,0,0)\}$ tendrá 2 componentes conexas. Por otro lado, cualquier abierto $U\subset \bb{R}^2$ al quitarle un punto $p\in U$ tendrá una única componente conexa por lo que no podrá existir ningún homeomorfismo entre $U$ y $V$.\\

    Para ver que $U\setminus\{p\}$ solo tiene una componente conexa tenemos que ver que al ser $U$ abierto existirá un abierto básico $A_p$ para todo $p\in U$ tal que $p\in A_p\subset U$. Además $A_p=B(p,\veps)$ con $\veps>0$ por lo que $A_p\setminus\{p\}$ seguirá siendo conexo. En consecuencia $U\setminus\{p\}$ será conexo.
\end{ejercicio}

\section{Aplicaciones diferenciables}

% PEDIR A DANI

\begin{definicion}
    Sean $S,S'\subset \bb{R}^3$ superficies y $O\subset \bb{R}^n$ un abierto. Dada $f$ una aplicación definida de alguna de estas formas:
    \begin{enumerate}
        \item[a)] $f:S\to \bb{R}^m$ con $m\geq 1$.
        \item[b)] $f:O\subset \bb{R}^n \to S$.
        \item[c)] $f:S \to S'$
    \end{enumerate}
    Diremos que $f$ es \textbf{diferenciable} si verifica respectivamente según su definición:
    \begin{enumerate}
        \item[a)] $f \circ X$ es diferenciable en el sentido de análisis para toda parametrización $X:U\to S$ (se tendrá $f \circ X: U\to \bb{R}^m$).
        \item[b)] $i \circ f: O\to \bb{R}^3$ es diferenciable en el sentido de análisis.
        \item[c)] $i' \circ f$ es diferenciable en el sentido de a)
    \end{enumerate}
    donde $i:S \to \bb{R}^3$, $i':S' \to \bb{R}^3$ denotan la inclusión.
\end{definicion}

\begin{propiedades}\
    \begin{enumerate}
        \item \textbf{Topológica:} Si $f$ es diferenciable en sentido a),b) o c), entonces $f$ es continua. 
        \begin{proof}
            Demostración
        \end{proof}
        \item \textbf{Algebraicas:} Sean $f,g:S\to \bb{R}^m$ con $m\geq 1$, $\lambda: S\to \bb{R}$ diferenciables (en el sentido a)). Entonces tendremos
        \begin{itemize}
            \item $f+g, \lambda f: S\to \bb{R}^m$ son diferenciables.
            \item $\langle f,g \rangle: S\to \bb{R}$ es diferenciable
            \item Además, si $m=3$ se tiene que $f \wedge g:S \to \bb{R}^3$ es diferenciable.
        \end{itemize}
        \begin{proof}
            Demostración
        \end{proof}
    \end{enumerate}
\end{propiedades}

\begin{ejemplo}\
    \begin{enumerate}
        \item Restricciones
        \item Constantes
        \item[$a_2$)] Distancias. Tomamos $p_0\in \bb{R}^3$, $O=\bb{R}^3$ y $F:\bb{R}^3 \to \bb{R}$ dada por 
        \begin{gather*}
            F(p) = |p-p_0|^2
        \end{gather*}
        que es la distancia al cuadrado en el punto $p_0$. Tenemos que $F$ es diferenciable en $O=\bb{R}^3$. Si tomamos la restricción $F_{|_S}:S \to \bb{R}$ sabemos por lo anterior que es diferenciable.\\

        Si definimos ahora $F(p)=|p-p_0|$ con $p_0\in \bb{R}^3$ tendremos que es diferenciable en $O = \bb{R}^3\setminus \{p_0\}$ por lo que necesitaremos $S\subset O \subset \bb{R}^3 \setminus \{p_0\}$, es decir, $p_0\notin S$. En este caso tendremos $F_{|_S}=f:S\to \bb{R}$ dada por $f(p)=|p-p_0|$ $\forall p \in S$ es diferenciable.\\

        Si tenemos $R\subset \bb{R}^3$ una recta, $p_0\in R\subset \bb{R}^3$ y $v\in \bb{R}^3\setminus\{0\}$ con $|v|=1$. Dado un punto $p\in \bb{R}^3$ queremos calcular la distancia de $p$ a $R$. Tendremos
        \begin{gather*}
            dist(p,R)^2 = |p-p_0| - \langle p-p_0, v \rangle^2
        \end{gather*}\\

        Sea $P\subset \bb{R}^3$ el plano que pasa por $p_0\in \bb{R}^3$ y tiene un vector unitario normal $v\in \bb{R}^3$ con $|v|=1$. Consideramos 
        \Func{F}{O=\bb{R}^3}{\bb{R}}{p}{F(p)=\langle p-p_0,v \rangle}
        que será la función \textbf{distancia orientada a $P$}. Tenemos que $F$ es diferenciable en el sentido de análisis. Consideramos $S\subset \bb{R}^3$ dada por $p\in S \mapsto \langle p-p_0, v\rangle$ que será la función \textbf{distancia orientada de los puntos de la superficie $S$ al plano $P$}.

        \item[a)] Es fácil o difícil de determinar si una función es diferenciable. Es claramente difícil porque intervienen \underline{todas} las parametrizaciones de la superficie. Consideramos el siguiente resultado:\\
        
        Dada $S\subset \bb{R}^3$ una superficie consideramos $X_i:U_i\to S$ con $i=1,2$ dos parametrizaciones de $S$ tales que 

        Suponemos que $X_1(U_1)\cap X_2(U_2) \neq \emptyset$
        \begin{gather*}
            \xymatrix{
                X_1^{-1} (X_1(U_1)\cap X_2(U_2)) \ar[r]^{X_2^{-1}\circ X_1} & X_2^{-1} (X_1(U_1)\cap X_2(U_2))
            }
        \end{gather*}
        Esta composición es un homeomorfismo entre abiertos de $\bb{R}^2$. Esta aplicación se llama \textbf{cambio de parámetros} o \textbf{cambio de coordenadas} (estamos haciendo 2 representaciones del mismo trozo de mapa). \\

        \textbf{Pequeño lema técnico}\footnote{``técnico'' significa que no lo explica (del griego \textit{techne})}\textbf{:} El cambio de parámetros es un difeomorfismo entre abiertos de $\bb{R}^2$.\\

        Una consecuencia de esto es la siguiente: Sea $f:S\to \bb{R}^m$ donde $S \subset \bb{R}^3$ es una superficie. Supongamos que $\forall p \in S$ existe una parametrización $X_p:U_p\to S$ talAplicaciones diferenciables que $p\in X_p(U_p)$ y $f\circ X_p$ es diferenciable en el sentido de análisis, entonces $f$ es diferenciable en el sentido a).
        \begin{proof}
            Sea $X:U\to S$ una parametrización arbitraria. Hay que ver que $f\circ X$ es diferenciable (en el sentido de análisis). Tenemos que $f\circ X : U \to \bb{R}^m$. Tomamos un $q\in U$ y le aplicamos la parametrización denotando $X(q)=p\in S$. Para dicho $p\in S$ existe por hipótesis una parametrización $X_p:U_p\to S$ tal que $p\in X_p(U_p)$ y tal que $f\circ X_p$ es diferenciable (AN\footnote{Usaremos esta notación para referirnos a ``en el sentido del análisis''.}). Para poder aplicar un cambio de cartas entre $X$ y $X_p$ necesitamos que se solapen las imágenes pero sabemos que 
            \begin{gather*}
                p\in (X_p(U_p)\cap X(U))=O \neq \emptyset
            \end{gather*}
            Tenemos entonces que $X_p^{-1}\circ X: X^{-1}(O)\to X_p^{-1}(O)$ es un difeomorfismo y además
            \begin{gather*}
                f\circ X_{|_{X^{-1}(O)}} = f_{|_O}\circ X_{|_{X^{-1}(O)}} = (f_{|_O}\circ X_p) \circ (X_p^{-1} \circ X_{|_{X^{-1}(O)}})
            \end{gather*}
            Sabíamos ya que $(X_p^{-1} \circ X_{|_{X^{-1}(O)}})$ es difeomorfismo (AN) y además, por hipótesis tenemos que $(f_{|_O}\circ X_p)$ es diferenciable (AN), por tanto llegamos a que $f\circ X_{|_{X^{-1}(O)}}$ es diferenciable. Además, como $q\in X^{-1}(O)$ tendremos que $f\circ X$ es diferenciable en $q$ para todo $q\in U$. Podemos concluir que $f\circ X$ es diferenciable (AN) para toda parametrización $X$.
        \end{proof}

        Una consecuencia de esto es que solo necesitaremos demostrar la diferenciabilidad en el sentido a) para un recubrimiento, en lugar de para toda parametrización.

        \item Supongamos que $S\subset \bb{R}^3$ es una superficie y $X:U\to S$ una parametrización. Queremos ver si las siguientes aplicaciones son diferenciables:
        \begin{enumerate}
            \item $X:U\to X(U)$. Es diferenciable en sentido b). Como $U\subset \bb{R}^2$ tendremos que $X(U)\subset \bb{R}^3$ superficie es diferenciable (AN). Por tanto sí es diferenciable.
            \item $X^{-1}:X(U)\to U$. Tenemos que esta aplicación sale de una superficie de $\bb{R}^3$ (al igual que en el apartado anterior) y llega a un abierto de un euclídeo por lo que estamos en la situación a). En este caso tenemos que $\forall p \in X(U)$ sirve como $X_p=X$, ya que $X$ es una parametrización que recubre a $X(U)$. Componemos entonces $X_p\circ f = X^{-1}\circ X = 1_U$ que es diferenciable (AN) por lo que concluimos que $X^{-1}$ sí es diferenciable.
        \end{enumerate}     
    \end{enumerate}
\end{ejemplo}

\subsection{Propiedades y ejemplos suplementarios}
\begin{enumerate}
    \item Una \textbf{curva diferenciable en una superficie} $S$ es una curva $\alpha:I\to \bb{R}^3$ diferenciable tal que su traza $\alpha(I)\subset S$. 
    
    \item Consideramos $S_1, S_2, S_3$  y $f_1, f_2$ dos aplicaciones diferenciables tales que 
    \begin{gather*}
        \xymatrix{
            S_1 \ar[r]^{f_1} & S_2 \ar[r]^{f_2} & S_3
        }
    \end{gather*}
    Si $S_1, S_2, S_3$ son superficies o abiertos de un euclídeo, entonces $f_2\circ f_1: S_1 \to S_3$ es diferenciable.

    \item Sea $S\subset \bb{R}^3$ una superficie y $O\subset S$ un abierto suyo. Supongamos que $f:S\to S'$ es diferenciable donde $S'$ es una superficie o un abierto de un euclídeo. Entonces $f_{|_O}:O\to S'$ es diferenciable.
    
    \item Consideramos $f:S\to S'$ donde $S\subset \bb{R}^3$ es una superficie y $S'$ es una superficie o un abierto de un euclídeo. Supongamos que podemos escribir
    \begin{gather*}
        S=\bigcup\limits_{i\in I} S_i\ \ \ \ S_i\subset S \text{ abierto }\ \ \ f_{|_{S_i}}:S_i\to S\ \ \text{ diferenciable}
    \end{gather*}
    Entonces $f:S\to S'$ también es diferenciable.

    \item Supongamos que $S,S'\subset \bb{R}^3$ son superficies y consideramos $f:S\to S'$. Si $f$ es diferenciable en el sentido c), es biyectiva, y su inversa $f^{-1}:S'\to S$ es diferenciable, diremos que $f$ es un \textbf{difeomorfismo}. \\
    
    La identidad de una superficie en sí misma y la composición de difeomorfismos también son difeomorfismos. Por eso $S$ es difeomorfa a $S'$ es una relación de equivalencia en la colección de todas las superficies.
\end{enumerate}

\begin{ejercicio}
    Demostrar que las superficies 
    \begin{gather*}
        C = \{(x,y,z)\in \bb{R}^3 : x^2 + y^2 = 1\}\\
        H = \{(x,y,z)\in \bb{R}^3 : x^2 + y^2 -z^2 = 1\}
    \end{gather*}
    son difeomorfas.\\

    En primer lugar veamos que son superficies. Como son cuádricas sin puntos singulares esto lo tenemos automáticamente. Veamos ahora que son difeomorfas. Construímos la siguiente aplicación $f : C \to H$:
    \begin{gather*}
        f(x,y,z) = (x\sqrt{1+z^2}, y\sqrt{1+z^2}, z)\in H
    \end{gather*}
    Veamos que $f(C)\subset H$, para lo cual tendremos que ver que cumple la ecuación de $H$:
    \begin{gather*}
        (x\sqrt{1+z^2})^2 + (y\sqrt{1+z^2})^2 - z^2 = (1+z^2)(x^2 + y^2) - z^2 = 1+z^2 -z^2 = 1
    \end{gather*}
    Vemos que es diferenciable en el sentido c) ya que si lo vemos como $f:C\to \bb{R}^3$ tendremos que ver que lo es en el sentido a). Además podemos ver $f = F_{|C}$ donde 
    \Func{F}{\bb{R}^3}{\bb{R}^3}{(x,y,z)}{(x\sqrt{1+z^2}, y\sqrt{1+z^2}, z)}
    por lo que $f$ es una restricción de una aplicación entre euclídeos por lo que solo tendremos que ver que es diferenciable en el sentido del análisis (efectivamente lo es). Tenemos entonces que $f$ es diferenciable. Consideramos 
    \Func{g}{H}{C}{(x,y,z)}{\left(\dfrac{x}{\sqrt{1+z^2}}, \dfrac{y}{1+z^2}, z\right)}
    Se deja como ejercicio ver que:
    \begin{itemize}
        \item Bien definida en $H$.
        \item Bien definida en $\bb{R}^3$.
        \item El mismo argumento para $f$ funciona para $g$.
        \item $f$ y $g$ son inversas.
    \end{itemize}
\end{ejercicio}

\section{Vector tangente a una Superficie}

\begin{definicion}
    Consideramos $S\subset \bb{R}^3$ una superficie, $p\in S$ y un vector $v\in \bb{R}^3$. Diremos que $v$ es \textbf{tangente} a $S$ en $p$ si existe una curva diferenciable $\alpha : (-\veps, \veps) \to S$ con $\veps > 0$ tal que
    \begin{gather*}
        \alpha(0)=p\ \ \ \text{ y } \ \ \ \alpha'(0)=v
    \end{gather*}
    Al \underline{conjunto} de vectores tangentes a la superficie $S$ en el punto $p$ lo representaremos como $T_pS \subset \bb{R}^3$.
\end{definicion}

\begin{observacion}
    Consideramos una superficie $S\subset \bb{R}^3$, $p\in S$ y $X:U\to S$ una parametrización de $S$ tal que $p\in X(U)$. En este caso tendríamos $X^{-1}(p)\in U \subset \bb{R}^2$. Consideramos $w\in \bb{R}^2$ arbitrario y podemos ver $t \mapsto X^{-1}(p)+tw$ para todo $t\in \bb{R}$, es decir, consideramos la aplicación $\beta: \bb{R} \to \bb{R}^3$ dada por 
    \begin{gather*}
        \beta(t) = X^{-1}(p) + tw
    \end{gather*}
    que verificará
    \begin{itemize}
        \item $\beta(0)=X^{-1}(p)$.
        \item $\beta'(0) = w$
    \end{itemize}
    Pero $tr(\beta)$ no tiene por qué estar contenida en $U$, de hecho lo más frecuente es que $tr(\beta)\not\subset U$ pero lo que sí se dará es que como $\beta$ es continuo, $\beta(0)=X^{-1}(p)$ y $U$ es un entorno abierto de $X^{-1}(p)$ tendremos que existirá un $\veps >0$ tal que 
    \begin{gather*}
        \beta (-\veps, \veps) \subset U
    \end{gather*}
    Sea $\alpha = X \circ \beta : (-\veps, \veps) \to S$ tendremos que será diferenciable y verificará 
    \begin{gather*}
        \alpha(0)=X(\beta(0)) = X(X^{-1}(p)) = p\\
        T_pS \owns \alpha'(0) = (X \circ \beta)'(0) \overset{(\ast)}{=} (dX)_{X^{-1}(p)}(\beta'(0)) = (dX)_{X^{-1}(p)}(w)
    \end{gather*}
    donde en $(\ast)$ hemos aplicado la regla de la cadena del análisis. Hemos probado que $\forall w \in \bb{R}^2$ tendremos que
    \begin{gather*}
        (dX)_{X^{-1}(p)}(w)\in T_pS
    \end{gather*}
    como $(dX)_{X^{-1}(p)} : \bb{R}^2 \to \bb{R}^3$ tendremos que $(dX)_{X^{-1}(p)}(\bb{R}^2)\subset T_pS$. Lo ideal es que se diera la igualdad (se va a dar). \\
    
    Para verlo tomemos $v\in T_pS$, o equivalentemente, un $v$ tangente a $S$ en $p$, es decir, tal que
    \begin{gather*}
        \exists \veps > 0,\ \ \exists \alpha:(-\veps, \veps) \to S\text{ diferenciable con }\ \ \alpha(0)=p,\ \ \alpha'(0)=v
    \end{gather*}
    Tenemos $X:U\to S$ una parametrización con $p\in X(U)$. Sabemos que $\alpha(-\veps, \veps)$ no tiene por qué estar dentro de $X(U)$ pero por la continuidad de $\alpha$ y sus propiedades tendremos que existirá un $0 < \delta \leq \veps$ tal que 
    \begin{gather*}
        \alpha : (-\delta, \delta) \to X(U)\subset S
    \end{gather*}
    Definimos $\beta = X^{-1} \circ \alpha : (-\delta, \delta) \to U\subset \bb{R}^2$ que verificará
    \begin{gather*}
        \beta(0) = X^{-1}(\alpha(0)) = X^{-1}(p)
    \end{gather*}
    Podemos escribir además $\alpha = X \circ \beta$ (despejando, lo que nos ayudará a derivar). Sabemos que
    \begin{gather*}
        v = \alpha'(0) \overset{(\ast)}{=} (dX)_{X^{-1}(p)}(\beta'(0))
    \end{gather*}
    donde en $(\ast)$ hemos vuelto a aplicar la regla de la cadena del análisis. Como $\beta'(0)\in \bb{R}^2$ tendremos que 
    \begin{gather*}
        v \in (dX)_{X^{-1}(p)}(\bb{R}^2)
    \end{gather*}
    por lo que tendremos la igualdad buscada.\\

    Hemos demostrado que si tenemos una superficie $S\subset \bb{R}^3$, $p\in S$, $X:U\to S$ una parametrización con $p\in X(U)$, entonces 
    \begin{gather*}
        T_pS = (dX)_{X^{-1}(p)}(\bb{R}^2) = img(dX)_{X^{-1}(p)}
    \end{gather*}
    Esto nos dará las siguientes consecuencias:
    \begin{enumerate}
        \item Como $(dX)_{X^{-1}(p)}$ es lineal e inyectiva tenemos que $T_pS$ es un subespacio vectorial de $\bb{R}^3$ de dimensión 2, es decir, $T_pS$ es un plano vectorial de $\bb{R}^3$ para todo $p\in S$. Normalmente se representará gráficamente como $p + \langle T_pS \rangle$ (el plano afín).
        
        \item $(dX)_{X^{-1}(p)}(\bb{R}^2) = img(dX)_{X^{-1}(p)}$, es decir, no depende de la parametrización $X$ elegida.
        
        \item $T_pS = img(dX)_{X^{-1}(p)} = (dX)_{X^{-1}(p)}(\bb{R}^2)$. Elegimos $\{(1,0), (0,1)\}$ la base canónica de $\bb{R}^2$. Entonces tendremos que 
        \begin{gather*}
            \{(dX)_{X^{-1}(p)}(1,0),\ (dX)_{X^{-1}(p)}(0,1)\} \text{ base de } T_pS
        \end{gather*}
        o equivalentemente
        \begin{gather*}
            \{X_u(X^{-1}(p)), X_v(X^{-1}(p))\} \text{ base de } T_pS
        \end{gather*}
        para toda parametrización $X$.
    \end{enumerate}
\end{observacion}

\begin{ejemplo}\
    \begin{enumerate}
        \item Sea $O\subset \bb{R}^3$ un abierto, $f:O\to \bb{R}$ diferenciable (AN), $a$ un valor regular de $f$. Suponemos $\emptyset \neq f^{-1}(\{a\}) = S$ y $p\in S$. Tenemos 
        \begin{gather*}
            \alpha : (-\veps, \veps) \to S\text{ con } \alpha(0)=p\ \ \text{ y } \ \ \alpha'(0)=v\in T_pS
        \end{gather*}
        Además $\alpha(t)\in S$ para todo $|t|<\veps$ por lo que $f(\alpha(t)) = a$ para todo $|t| < \veps$. Tenemos que 
        \begin{gather*}
            (df)_p(v) = 0 \Rightarrow v \in \ker(df)_p\ \ \ \forall v \in T_pS
        \end{gather*}
        Tenemos que $T_pS \subset \ker(df)_p$ y como ambos espacios tienen dimensión 2 se dará la igualdad, es decir,
        \begin{gather*}
            T_pS = \ker(df)_p = \langle \nabla f_p \rangle^{\bot}
        \end{gather*}

        \begin{enumerate}
            \item Consideramos un plano $P$, un punto $p_0\in P$ y $a$ el vector normal unitario de $P$. De esta forma podemos escribir:
            \begin{gather*}
                P = \{p\in \bb{R}^3 : \langle p - p_0, a\rangle = 0\} = f^{-1}(\{0\})
            \end{gather*}
            y definimos 
            \Func{f}{\bb{R}^3}{\bb{R}}{p}{f(p)=\langle p -p_0, a\rangle}
            Se deja como ejercicio comprobar que $0$ es valor regular de $f$.
            De esta forma, dado un $p\in P$ tenemos que 
            \begin{gather*}
                T_pP = \langle (\nabla f)_p \rangle^{\bot} = \ker(df)_p
            \end{gather*}
            donde
            \Func{(df)_p}{\bb{R}^3}{\bb{R}}{v}{(df)_p(v)}
            Recordemos que $(df)_p(v)=(DF)(p,v)$ (con la notación del análisis\footnote{Es la última vez que lo pone de verdad, lo jura por Snoopy}). Calculemos dicha diferencial:
            \begin{align*}
                (df)_p(v)&= \frac{d}{dt}|_{t=0} f(p+tv) = \frac{d}{dt}|_{t=0} \langle p + tv - p_0, a \rangle = \langle v,a \rangle\\
                (\nabla f)_p &= a \ \ \ \forall p \in P
            \end{align*}
            De esta forma podemos asegurar que 
            \begin{gather*}
                T_pP = \langle a \rangle^{\bot}
            \end{gather*}

            \item \begin{definicion}
                Dada una superficie $S\subset \bb{R}^3$ y un punto $p\in S$ llamamos \textbf{recta normal} a $S$ en $p$ a la recta afín que pasa por $p$ y es perpendicular a $T_pS$.
            \end{definicion}

            De esta forma hemos probado que todas las rectas normales a un plano son paralelas. Esto deja abierta la pregunta de si una superficie con todas sus rectas normales paralelas es necesariamente un plano (o un trozo abierto de plano). Probablemente esto sea verdad.
        \end{enumerate}
        \item Dado $p_0\in \bb{R}^3$, $R>0$ y consideramos la superficie $S$ dada por 
        \begin{gather*}
            S = \{p\in \bb{R}^3 : |p-p_0|^2 = R^2\} = S^2(p_0, S)
        \end{gather*}
        Podemos definir la aplicación 
        \Func{g}{\bb{R}^3}{\bb{R}}{p}{g(p)=|p-p_0|^2}
        y podemos así ver la superficie como $S=g^{-1}(\{R^2\})$. Ya estudiamos que $R^2$ es un valor regular de $g$ (una función diferenciable). Tenemos entonces que $S$ es una superficie y acabamos de ponerla en forma implícita. Sabemos que 
        \begin{gather*}
            T_pS^2(p_0, R) = \ker(dg)_p = \langle \nabla g_p \rangle^{\bot}
        \end{gather*}
        consideramos la diferencial:
        \Func{(dg)_p}{\bb{R}^3}{\bb{R}}{v}{(dg)_p(v)=2\langle p -p_0, v\rangle}
        Con lo que tendremos 
        \begin{gather*}
            (\nabla g)_p = 2(p-p_0) \Rightarrow T_pS^2(p_0, S) = \langle p-p_0\rangle ^{\bot}
        \end{gather*}
        Nos planteamos ahora cómo son las rectas normales a la esfera $S^2(p_0, S)$ en un punto $p\in S^2(p_0, S)$. Tendrán que pasar por $p$ y llevar la dirección $p-p_0$ por lo que podríamos escribir la ecuación paramétrica
        \begin{gather*}
            p+ \lambda (p-p_0) \ \ \ \lambda \in \bb{R}
        \end{gather*}
        Tenemos por tanto que $\forall p \in S^2(p_0, R)$ existe un $\lambda(p)=-1$ tal que $p-(p-p_0) = p_0$ lo que nos dice que todas las rectas normales a una esfera son \textbf{concurrentes} (todas pasan por un mismo punto, el origen). \\

        De nuevo se plantea la cuestión de que si todas las rectas normales a una superficie son concurrentes, entonces necesariamente será un trozo abierto de esfera (sospechamos que sí).

        \item Consideramos ahora $S\subset \bb{R}^3$ una superficie, un subconjunto abierto $S'\subset S \subset \bb{R}^3$ y un punto $p'\in S'\subset S$. Sabemos que $S'$ es una superficie y queremos comparar $T_{p'}S'$ y $T_pS$. Sabemos que para cada $v\in T_{p'}S'$ existe una aplicación $\alpha: (-\veps, \veps) \to S'\subset S$ con 
        \begin{gather*}
            \alpha(0) = p'\ \ \text{ y } \ \ \alpha'(0) = v
        \end{gather*}
        lo que nos dice que $v\in T_pS$ por lo que tendremos que $T_{p'}S'\subset T_pS$ lo que finalmente (con un razonamiento simétrico) nos dará la igualdad:
        \begin{gather*}
            T_{p'}S' = T_pS
        \end{gather*}

        \begin{enumerate}
            \item Consideramos la aplicación $f:S\to \bb{R}^m$ diferenciable. Sabemos que para cada $v\in T_pS$ existe una aplicación $\alpha: (-\veps, \veps) \to S$ tal que $\alpha(0)=p$ y $\alpha'(0)=v$.  Definimos así $(df)_p$ como 
            \Func{(df)_p}{T_pS}{\bb{R}^m}{v}{(df)_p(v) = \frac{d}{dt}|_{t=0}(f(\alpha(t))) = (f\circ \alpha)'(0)}
            Además, para que tenga sentido la expresión de la derecha necesariamente necesitaremos $v\in T_pS$.\\

            Hay que ver que esta función está bien definida, en concreto ver que $(f\circ \alpha)'(0)$ no depende de $\alpha$, sino de $\alpha'(0)=v$.\\

            Consideramos una parametrización $X:U\to S$ y un punto $p\in X(u)$. Podemos considerar por continuidad que $\alpha(t)\in X(U)$ para todo $t\in (-\veps, \veps)$. Podemos escribir
            \begin{gather*}
                f\circ \alpha = (f\circ X)\circ (X^{-1}\circ \alpha)
            \end{gather*}
            Como ambas composiciones son difereciables (AN) podemos derivar en $0$ y por la regla de la cadena:
            \begin{gather*}
                (f\circ \alpha)'(0) = d(f\circ X)_{X^{-1}(p)}(X^{-1}\circ \alpha)'(0)
            \end{gather*}
            Sabemos que $X\circ X^{-1}\circ \alpha = \alpha$ y 
            \begin{gather*}
                v = \alpha'(0) = (dX)_{X^{-1}(p)}(X^{-1}\circ \alpha)'(0)
            \end{gather*}
            por lo que despejando
            \begin{gather*}
                (X^{-1}\circ \alpha)'(0) = [(dX)_{X^{-1}(p)}]^{-1}(v)
            \end{gather*}
            y tenemos finalmente que
            \begin{gather*}
                (df)_p(v) = d(f\circ X)_{X^{-1}(p)}[(dX)_{X^{-1}(p)}]^{-1}(v)
            \end{gather*}
            y llegamos a que no depende de $\alpha$ sino solo de $v$ y además lo hace linealmente.

            \item Consideramos $f:O\subset \bb{R}^n \to S$ con $O$ abierto y un punto $x\in O$. Si repetimos la definición tendremos 
            \begin{gather*}
                (df)_X ; \bb{R}^n \to T_{f(x)S\subset \bb{R}^3}
            \end{gather*}

            \item Supongamos ahora $f:S\to S'$, $p\in S$. Si repetimos la definición 
            \begin{gather*}
                (df)_p : T_pS \to T_{f(p)}S'
            \end{gather*}
        \end{enumerate}
    \end{enumerate}
\end{ejemplo}

\begin{ejemplo}\
    \begin{enumerate}
        \item Restricciones. Consideramos $F:O\subset \bb{R}^3$ diferenciable (AN) y $S\subset O$. Si llamamos $f=F_{|S}:S\to \bb{R}^m$ sabemos ya que es diferenciable por lo que podemos tomar $p\in S$ y existe la diferencial que será una aplicación $(df)_p: T_pS\to \bb{R}^m$. Aplicando la definición tendremos
        \begin{gather*}
            (df)_p(v)=(f \circ \alpha)'(0)
        \end{gather*}
        donde $\alpha:(-\veps, \veps) \to S$ con $\alpha(0)=p$ y $\alpha'(0)=v$. Además en este caso $f$ y $F$ actúan igual en el dominio por lo que 
        \begin{gather*}
            (df)_p(v)=(F \circ \alpha)'(0) = (dF)_p(v)
        \end{gather*}
        y esto ocurre para todo $v\in T_pS$. Esto nos dice que ante restricciones las diferenciales actúan iguales pero no son iguales el todo ya que $(df)_p : T_pS \to \bb{R}^m$ y $(dF)_p : \bb{R}^3 \to \bb{R}^m$ por lo que concluímos que 
        \begin{gather*}
            (dF_{|S})_p = (dF)_{|T_pS}
        \end{gather*}
        es decir, la diferencial de la restricción es la restricción de la diferencial en el plano tangente.
        
        \item Teníamos la aplicación $F:\bb{R}^3\to \bb{R}$ dada por $F(p)=\langle p -p_0, u\rangle$ con $|u|^2=1$, $p_0\in \bb{R}^3$. Notamos $h=F_{|S}$. Queremos calcular $(dh)_p(v)$ y aplicando lo recién visto para restricciones
        \begin{gather*}
            (dh)_p(v) = (dF)_p(v) = \langle v,u \rangle
        \end{gather*}

        \item Consideramos ahora la distancia a un punto, es decir, $F:\bb{R}^3\to \bb{R}$ dada por $F(p)=|p-p_0|^2$. Notamos $f=F_{|S}$ con $S\subset \bb{R}^3$.
        \begin{gather*}
            (df)_p(v) = (dF)_p(v) = 2\langle p-p_0, v \rangle
        \end{gather*}

        \item Definimos $F:\bb{R}^3 \setminus \{p_0\} \to \bb{R}$ dada por $F(p)=|p-p_0|$. Sabemos que $p_0\notin S$ y además, notando $f= F_{|S}:S\to \bb{R}$ tenemos 
        \begin{gather*}
            (df)_p(v)=\frac{2\langle p-p_0,v\rangle}{2|p-p_0|} = \frac{\langle p-p_0,v\rangle}{|p-p_0|}
        \end{gather*}
        de esta forma tendremos también
        \begin{gather*}
            (d\ dist^2_{|S})_p(v) = 2\langle p-p_0,v\rangle - 2 \langle p-p_0, u\rangle \langle v,u\rangle
        \end{gather*}

        \item Inclusiones. Consideramos una superficie $S\subset \bb{R}^3$ y la aplicación inclusión $i:S \to \bb{R}^3$. Dado un punto $p\in S$ podemos considerar $(df)_p : T_pS \to \bb{R}^3$. Podemos considerar $i=I_{3|_S}=i_S$ y por tanto
        \begin{gather*}
            (di)_p = (dI_{3|_S})_p = (dI_3)_p = I_3 \Rightarrow (di)_p = I_{3|_{T_pS}} = i_{T_pS}
        \end{gather*}

        \item Aplicaciones constantes. Consideramos la superficie $S\subset \bb{R}^3$ y la aplicación $f:S\to S'$ (superficie o abierto de $\bb{R}^n$). Si $f$ es constante entonces existe un $a\in S'$ tal que $f(p)=a$ para todo $p\in S$. De esta forma, para un cierto $p\in S$ tendremos que $(df)_p : T_pS \to T_aS'$ donde aplicando la definición existirá un $\alpha:(-\veps, \veps) \to S$ con $\alpha(0)=p$ y $\alpha'(0)=v$ de forma que
        \begin{gather*}
            (df)_p(v) = (f\circ \alpha)'(0) = 0 \ \ \ \forall v\in T_pS
        \end{gather*}
        ya que $\forall t \in (-\veps, \veps)$ tendremos que $(f\circ \alpha)(t) = f(\alpha(t)) = a$ por lo que $f\circ \alpha$ será constante. Por tanto 
        \begin{gather*}
            f:S \to S' \ \text{ es constante } \Rightarrow (df)_p = 0 \ \ \ \forall p \in S
        \end{gather*}
    \end{enumerate}
\end{ejemplo}

\begin{propiedades}
    Algunas propiedades importantes que podemos destacar son las siguientes:
    \begin{enumerate}
        \item \textbf{Regla de la cadena.} Si estamos en la siguiente situación
        \begin{gather*}
            \xymatrix{
                S \ar[r]^f & S' \ar[r]^g & S''
            }
        \end{gather*}
        donde $S,S',S''\subset \bb{R}^3$ son superficies o abiertos de un espacio euclídeo y $f,g$ son aplicaciones diferenciables. Entonces $\forall p \in S$ se tiene que 
        \begin{gather*}
            d(g\circ f)_p = (dg)_{f(p)}\circ (df)_p
        \end{gather*}
        \begin{proof}
            Tenemos que pensar en el siguiente diagrama
            \begin{gather*}
                \xymatrix{
                    T_pS \ar[rr]^{d(g\circ f)_p} \ar[dr]_{(df)_p} && T_{g(f(p))}S''\\
                    & T_{f(p)}S' \ar[ur]_{(dg)_{f(p)}}
                }
            \end{gather*}
            y queremos ver que conmuta. Aplicando la definición tenemos que para todo $v\in T_pS$ existirá un $\alpha: (-\veps, \veps) \to S$ con $\alpha(0)=p$ y $\alpha'(0)=v$. Tendremos entonces 
            \begin{gather*}
                (f\circ \alpha):(-\veps,\veps) \ \ \text{ con }\ (f\circ \alpha)(0) = f(p) \ \ ,\ \ (f\circ \alpha)'(0) = (df)_p(v)
            \end{gather*}
            Si aplicamos ahora la definición de diferencial a $g$ tendremos 
            \begin{align*}
                d(g\circ f)_p(v) &= ((g\circ f)\circ \alpha)'(0) = (g\circ(f\circ \alpha))'(0) = (dg)_{f(p)}((df)_p(v))\\
                &= [(dg)_{f(p)} \circ (df)_p](v)\ \ \ \ \forall v \in T_pS
            \end{align*}
            Por lo que podemos concluir que 
            \begin{gather*}
                d(g\circ f)_p = (dg)_{f(p)} \circ (df)_p
            \end{gather*}
        \end{proof}

        \item \textbf{Extremos locales de una función.} Consideramos $S\subset \bb{R}^3$ una superficie y $f:S\to \bb{R}$ diferenciable. Suponemos que $f$ tiene un extremo local en $p_0\in S$. Entonces 
        \begin{gather*}
            (df)_{p_0} = 0
        \end{gather*}
        \begin{proof}
            Aplicando la definición a $(df)_{p_0} : T_{p_0}S \to \bb{R}$ tenemos que para $v\in T_{p_0}S$ existe un $\alpha:(-\veps, \veps) \to S$ con $\alpha(0)=p_0$ y $\alpha'(0)=v$. Además,
            \begin{gather*}
                (df)_{p_0}(v) = (f\circ \alpha)'(0)
            \end{gather*}
            Si $p_0$ es extremo local (por ejemplo un máximo local, sin pérdida de generalidad), entonces tenemos que $\exists V$ entorno abierto de $p_0$ en $S$ de forma que 
            \begin{gather*}
                f(p) \leq f(p_0)\ \ \ \ \forall p \in V
            \end{gather*}
            Podemos suponer que $\alpha(t)\in V$ para todo $|t|<\veps$ y entonces
            \begin{gather*}
                (f\circ \alpha)(t)= f(\alpha(t)) \leq f(p_0)=(f\circ \alpha)(0)
            \end{gather*}
            por lo que tenemos $f\circ \alpha : \bb{R}\to \bb{R}$ que tiene un máximo local en $t=0$. Del análisis sabemos que su derivada ha de anularse luego
            \begin{gather*}
                (f\circ \alpha)'(0)=0 \Rightarrow (df)_{p_0}(v)=0 \ \ \ \forall v \in T_{p_0}S
            \end{gather*}
        \end{proof}

        \item[2.bis.] Sea $S\subset \bb{R}^3$ una superficie conexa y $f:S\to S'$ (superficie o abierto de un euclídeo) diferenciable. Si $(df)_p=0$ para todo $p\in S$, entonces $f$ es constante.
        
        \begin{proof}
            Dado $a\in Im(f)$ podemos definir
            \begin{gather*}
                A=\{p\in S : f(p)=a\} = f^{-1}(\{a\}) \subset S
            \end{gather*}
            Sabemos que $A\neq \emptyset$ porque $a\in Im(f)$ y además es cerrado por ser imagen inversa de un cerrado en $\bb{R}$. Basta ver que $A$ es abierto en $S$. Dado un $p\in A$ podemos considerar una parametrización $X:U \subset \bb{R}^2 \to S$ tal que $p\in X(U)$. Tenemos que 
            \begin{gather*}
                f_{|X(U)} = (f\circ X) \circ X^{-1}
            \end{gather*}
            Aplicando la regla de la cadena recién probada tenemos que 
            \begin{gather*}
                0 = (df)_q = d((f\circ X)\circ X^{-1})_q = d(f\circ X)_{X^{-1}(q)} \circ (dX^{-1})_q
            \end{gather*}
            Como $(dX^{-1})_q$ es un isomorfismo, necesariamente tenemos que 
            \begin{gather*}
                d(f\circ X)_{X^{-1}(q)}=0
            \end{gather*}
            Aplicando propiedades de análisis sabemos que $f\circ X : U\subset \bb{R}^2 \to \bb{R}^m$ con diferencial cero en cada punto por lo que $f\circ X$ es constante lo que nos dice que 
            \begin{gather*}
                f_{|X(U)} = a \ \ \ \text{ constante} \Rightarrow p\in X(U)\subset A
            \end{gather*}
            Finalmente hemos probado que $A$ contiene un entorno abierto para cada $p\in A$, luego $A$ es abierto.
        \end{proof}

        \item \textbf{Teorema de la función inversa.} Consideramos $f:S\to S'$ donde $S,S'$ son superficies o abiertos de $\bb{R}^2$ diferenciable. Sea $p_0\in S$ donde $T_{p_0}S \to T_{f(p_0)}S'$ es regular (isomorfismo, rango 2, inyectiva). Entonces existe un entorno abierto $V$ de $p_0$ en $S$ y un entorno abierto $W$ de $f(p_0)$ en $S'$ tales que 
        \begin{gather*}
            f(V)=W \ \ \ \ f_{|V}: V \to W \ \ \text{ es un difeomorfismo}
        \end{gather*}
        Consecuencias
        \begin{enumerate}
            \item Si $(df)_p$ es regular para todo $p\in S$ entonces $f$ es un difeomorfismo local. Por tanto, $f$ es abierta y así $f(S)$ es abierto en $S'$.
            
            \begin{ejemplo}
                Sean $S'\overset{i}{\subset}S \subset \bb{R}^3$ superficies, $p'\in S'$ con $(di)_{p'} = \restr{i}{T_{p'}S'}$ donde $(di)_p'(T_{p'}S')\subset T_{p'}S$ y además tienen la misma dimensión, luego $(di)_p'$ es regular para todo $p'\in S'$. Si aplicamos la consecuencia recién vista, en particular $i(S')=S'\subset S$ es abierto. 
            \end{ejemplo}
        \end{enumerate}
    \end{enumerate}
\end{propiedades}

\begin{ejercicio}\
    \begin{enumerate}
        \item Consideramos $A\in M_3(\bb{R})$ con $A=A^t$. Definimos 
        \Func{f}{S^2}{\bb{R}}{p}{f(p)=\langle Ap, p \rangle}
        que es diferenciable por ser restricción de una aplicación diferenciable ne $\bb{R}^3$ por lo que en particular es continua. Topológicamente además sabemos que $S^2$ es compacta. Sabemos que $f$ alcanza su máximo en algún punto $p_M\in S^2$ y su mínimo en un punto $p_m\in S^2$. Hagamos ahora una distinción de casos:
        \begin{enumerate}
            \item Si $f(p_M)=f(p_m)$, entonces $f$ es constante y tendremos para cada $p\in S^2$:
            \begin{gather*}
                0=(df)_p : T_pS^2 \to \bb{R}
            \end{gather*}
            Ahora para cada $v\in T_pS^2$ tendremos que 
            \begin{gather*}
                0 = (df)_p(v) = \langle Av, p \rangle + \langle Ap, v \rangle = 2 \langle Ap, v \rangle
            \end{gather*}
            lo que nos dice que $\forall p \in S^2$ se tiene que
            \begin{gather*}
                \langle Ap, v \rangle = 0 \ \ \ \ \forall v \in T_pS^2 = \langle p \rangle^{\bot}
            \end{gather*}
            por lo que $Ap \bot v$ para todo $v\in \langle p \rangle^\bot$. Es decir, que para todo $p$ existirá un cierto $\lambda(p)$ tal que 
            \begin{gather*}
                Ap=\lambda(p)p
            \end{gather*}
            Además, $\lambda = cte = f(p) = \lambda(p)$ lo que significa que hemos podido diagonalizar la matriz (hemos obtenido los valores propios de $A$). Además, tendremos
            \begin{gather*}
                A = \lambda I_3 \ \ \ \ \lambda \in \bb{R}
            \end{gather*}
            y todas las bases ortonormales de $\bb{R}^3$ diagonalizan $A$ y sus valores propios son $\lambda, \lambda, \lambda$.

            \item Si $f(p_M)\neq f(p_m)$ (es claro que entonces $f(p_M)>f(p_m)$), entonces necesariamente $p_M\neq p_m$. Además, en particular $p_M$ y $p_m$ serán extremos locales (y de hecho globales) de $f$ por lo que, por las propiedades vistas anteriormente tendremos
            \begin{gather*}
                (df)_{p_M} = (df)_{p_m} = 0
            \end{gather*}
            Veamos lo que podemos deducir de esto:
            \begin{gather*}
                (df)_{p_M} = 0 \Rightarrow (df)_{p_M}(v) = 0 \ \ \forall v \in T_{p_M}S^2 
            \end{gather*}
            Siguiendo el razonamiento del caso anterior tendremos
            \begin{gather*}
                \langle Ap_M, v\rangle = 0 \ \ \forall v \in T_{p_M}S^2= \langle p_M \rangle^\bot
            \end{gather*}
            lo que nos dice que 
            \begin{gather*}
                Ap_M = \lambda p_M\ \ \text{ para cierto } \lambda\in \bb{R} \Rightarrow f(p_M)= \langle Ap_M, p_M \rangle = \lambda |p_M|^2 = \lambda
            \end{gather*}
            Tendremos entonces
            \begin{gather*}
                Ap_M = f(p_M)p_M
            \end{gather*}
            por lo que $p_M$ es un vector propio para $A$ con valor propio $f(p_M)$. \\
            
            Análogamente llegaríamos a que 
            \begin{gather*}
                Ap_m = f(p_m)p_m
            \end{gather*}
            por lo que $p_m$ es un vector propio para $A$ con valor propio $f(p_m)$.\\
            
            Veamos que forman una base ortonormal. Es claro que tienen módulo 1 ya que están en la esfera de radio 1. Veamos entonces que son ortogonales. Consideramos 
            \begin{align*}
                (f(p_M)-f(p_m))\langle p_M, p_m\rangle &= f(p_M)\langle p_M, p_m \rangle - f(p_m)\langle p_M, p_m \rangle = \\
                &= \langle f(p_M)p_M \rangle - \langle p_M, f(p_m)p_m\rangle = \\
                &= \langle Ap_M, p_m \rangle - \langle p_M, Ap_m \rangle =\\
                &= 0
            \end{align*}
            Ya que $\langle Ap_M, p_m \rangle=\langle p_M, Ap_m \rangle$ (la matriz es autoadjunta). Como teníamos $(f(p_M)-f(p_m))\neq 0$, necesariamente se dará $\langle p_M, p_m\rangle=0$, por lo que efectivamente son ortogonales. Por tanto, $\{p_M, p_m\}$ son dos vectores propios para $A$ unitarios y perpendiculares. Nos falta uno para tener una base ortonormal. Elegimos así un $p_0\in S^2$ tal que $p_0\bot p_M$ y $p_0\bot p_m$. Tendremos
            \begin{gather*}
                \langle Ap_0, p_M \rangle = \langle p_0, Ap_M \rangle = f(p_M) \langle p_0, p_M \rangle = 0\\
                \langle Ap_0, p_m \rangle = \langle p_0, Ap_m \rangle = f(p_m) \langle p_0, p_m \rangle = 0
            \end{gather*}
            luego el vector $Ap_0$ no tiene componente ni en $p_M$ ni en $p_m$, luego solo tendrá componente en $p_0$. Esto implica que 
            \begin{gather*}
                Ap_0 = f(p_0)
            \end{gather*}
            por lo que $p_0$ es propio para $A$ con valor propio $f(p_0)$. Tendremos finalmente que $\{p_M, p_m, p_0\}$ es una base ortonormal de $\bb{R}^3$ formada por vectores propios para $A$ con valores propios $\{f(p_M), f(p_m), f(p_0)\}$, base que diagonaliza a $A$.
        \end{enumerate}

        \item Dada $S\subset \bb{R}^3$ veamos que $S^\circ = \emptyset$. Veámoslo por reducción al absurdo. Si $p\in S^\circ$, entonces $\exists \veps > 0$ tal que $B(p,\veps)\subset S$. Tomemos $\{e_1, e_2, e_3\}$ la base canónica de $\bb{R}^3$ y consideramos para $i\in\{1,2,3\}$ las aplicaciones
        \Func{\alpha_i}{\bb{R}}{\bb{R}^3}{t}{\alpha_i(t)=p+te_i}
        Entonces si $|t|<\veps$ tendremos que $\alpha_i(t)\in B(p,\veps)\subset S$ (los vectores de la base son unitarios por lo que no se saldrán de la bola). Además tendremos que para cada $i\in \{1,2,3\}$
        \begin{gather*}
            \alpha_i(0)=p, \ \ \ \alpha_i'(0)=e_i\in T_pS
        \end{gather*}
        por lo que habremos encontrado tres vectores linealmente independientes en un plano, lo cual es imposible. Llegamos a contradicción.

        \item Consideramos $S\subset \bb{R}^3$ una superficie y un punto $p_0\in \bb{R}^3$. Se nos pide estudiar los puntos críticos de la función distancia al cuadrado a $p_0$ y extraer todas las consecuencias geométricas que se puedan.
        
        \begin{definicion}
            Un punto crítico $p\in S$ para la función $f:S\to \bb{R}$ es un punto donde $(df)_p = 0$.
        \end{definicion}

        Consideramos entonces $f(p)=|p-p_0|^2$ para todo $p\in S$. Tendremos entonces que $p$ es punto crítico si y solo si
        \begin{align*}
            (df)_p = 0 &\sii (df)_p(v) = 0 \ \ \ \forall v \in T_pS \sii
            \langle p-p_0, v \rangle = 0 \ \ \forall v\in T_pS \\
            &\sii p_0-p \bot v \ \ \forall v \in T_pS 
        \end{align*}
        lo cual ocurrirá si y solo si $p_0$ está en la recta normal a $S$ en $p$. Si $S$ es compacta, entonces hay al menos un punto crítico. Por $p_0$ se puede trazar una recta perpendicular a $S$.\\

        Si $S$ es conexa, entonces todas las rectas normales a $S$ son concurrentes (pasan por un mismo punto). Sea $p_0$ el punto común a todas las rectas normales. Consideramos la función distancia al cuadrado a $p_0$, $f$. Sabemos que $p_0$ está en la recta normal a $S$ en $p$ para todo $p\in S$. Además, $p$ es crítico para $f$ para todo $p\in S$ por lo que 
        \begin{gather*}
            (df)_p = 0 \ \ \ \forall p \in S
        \end{gather*}
        Esto nos dice que $f$ es constante, luego se verificará
        \begin{gather*}
            |p-p_0|^2 = R^2\ \ \ \forall p\in S,\ \ R \geq 0
        \end{gather*}
        Si $R=0$, tendríamos que $S=\{p_0\}$ pero esto no puede ser, por lo que $R>0$. Esto nos dice que $S\subset S^2(p_0, R)$, es decir, $S$ es un abierto de una esfera de $\bb{R}^3$.

        \item Buena pregunta (solo Dios sabe qué ejercicio va aquí).
        
        \item Consideramos $S\subset \bb{R}^3$ una superficie cerrada, un punto $a\in \bb{R}^3$ arbitrario. Demostrar que la distancia de $S$ al punto $a$ alcanza un mínimo.\\
        
        Consideramos el siguiente conjunto
        \begin{gather*}
            A=\{|p-a| : p\in S\}\subset \bb{R}
        \end{gather*}
        Sabemos que $A$ está acotado inferiormente por $0$. Por el axioma del supremo (del ínfimo) tendremos que
        \begin{gather*}
            \exists r = \inf_{p\in S} |p-a| \Rightarrow \exists \{p_n\}_{n\in \bb{N}}\subset S \text{ tal que } r=\lim_{n\to \infty} |p_n-a|
        \end{gather*}
        Podemos considerar la sucesión $\{|p_n-a|\}_{n\in \bb{N}}\subset \bb{R}$ que estará acotada. Por tanto existirá un $C\in \bb{R}$ tal que $|p_n-a|\leq C$ lo que nos dice que $\{p_n\}_{n\in \bb{N}}\subset B(a,R)$ para un radio $R>0$. Por tanto $\{p_n\}$ está acotada y por el teorema de Bolzano-Weierstrass existirá una parcial ${p_{n_k}}_{k\in \bb{N}}$ de $\{p_n\}_{n\in \bb{N}}$ convergente, es decir,
        \begin{gather*}
            \exists \lim_{k\to \infty} p_{n_k} = p \in \bar{S} = S
        \end{gather*}
        Donde en la última igualdad hemos aplicado que $S$ es cerrada. Por la unicidad del límite tendremos
        \begin{gather*}
            r = \lim_{n\to \infty} |p_n-a| = \lim_{k\to \infty} |p_{n_k}-a| \overset{(\ast)}{=} \left|\lim_{k\to \infty} p_{n_k}-a \right| = |p-a|,\ \ \ p\in S
        \end{gather*}
        Donde en $(\ast)$ hemos aplicado que la función distancia es continua. Tendremos entonces que $r$ es un mínimo en $S$ que se alcanza en $p$. Esto no funciona para el máximo.

        \item Sea $S\subset \bb{R}^3$ una superficie compacta. Probar que existe una recta $R$ de $\bb{R}^3$ que corta a $S$ perpendicularmente en al menos dos puntos suyos.\\
        
        Definimos la aplicación
        \Func{f}{S\times S}{\bb{R}}{(p,q)}{f(p,q)=|p-q|^2}
        que es continua. Como $S\times S$ es compacto tendremos que $f$ alcanza su máximo en algún punto $(a,b)\in S\times S$, lo que significa que 
        \begin{gather*}
            |p-q|^2 \leq |a-b|^2 \ \ \forall p,q\in S
        \end{gather*}
        Además, si $a=b$ significa que $|p-q|^2 = 0$ para todo $p,q\in S$, lo que implicaría que $S=\{p_0\}$ lo que nos llevaría a contradicción. Por tanto $a\neq b$ y podemos considerar $R$ la única recta de $\bb{R}^3$ que pasa por $a$ y $b$. Tenemos también
        \begin{gather*}
            |p-b| \leq |a-b|^2 \ \ \forall p \in S \Rightarrow dist^2(b) \ \ \text{ alcanza su máximo en }a\\
            |a-q| \leq |a-b|^2 \ \ \forall q \in S \Rightarrow dist^2(a) \ \ \text{ alcanza su máximo en }b
        \end{gather*}
        Por tanto $a$ está en la recta normal a $S$ que pasa por $b$ y $b$ está en la recta normal a $S$ que pasa por $a$ (tal y como lo hemos visto en ejemplos anteriores).
    \end{enumerate}
\end{ejercicio}